%% sections/dataset.tex
\section{Dataset}
\label{sec:dataset}

We present \textbf{KonkaniAMR}, the first Abstract Meaning Representation dataset for the Konkani language. The dataset contains 1,100 sentence–graph pairs and is split into training (880), validation (55), and test (165) sets.

\subsection{Data Collection}

\textbf{BPCC sample.} The primary source of sentences is the AI4Bharat BPCC (Bharat Parallel Corpus Collection) dataset \cite{gala2023indictrans2}, a large-scale parallel corpus of 22 Scheduled Indian languages paired with English, assembled from government documents, news, and publicly available web sources. We randomly sampled 1,000 Konkani sentences from the Konkani portion of BPCC. Sentences were filtered to exclude those shorter than four words or longer than 50 words, as very short sentences are semantically trivial and very long sentences place excessive demands on the annotator. The resulting sample spans a range of domains including news, administrative text, and factual descriptions, providing lexical and topical diversity.

\textbf{Newspaper sample.} To complement the BPCC material, we collected an additional 100 sentences from a Konkani newspaper. These sentences capture more colloquial registers and topic domains—sport, local politics, community events—that are underrepresented in the government-domain BPCC corpus. Including newspaper text ensures the dataset reflects the breadth of language use rather than a single genre.

\subsection{Automated Annotation with Gemini 2.5 Pro}

Manual AMR annotation is extremely labour-intensive: trained annotators require weeks of training and produce only a handful of annotations per hour. For a low-resource language like Konkani, where few annotators with both linguistic expertise and AMR knowledge are available, manual annotation of even a small dataset is prohibitively costly.

We therefore adopted a silver-annotation strategy, using Gemini 2.5 Pro as an automatic AMR annotator. For each sentence, we prompted the model with a structured instruction that described the AMR formalism, the Penman notation, and relevant PropBank predicate conventions. The model was asked to produce a well-formed, syntactically correct Penman graph that faithfully represents the semantic content of the Konkani sentence.

Each prompt was structured as follows:
\begin{enumerate}
    \item A system instruction explaining AMR, Penman notation, and key relation types (core arguments \texttt{:ARG0–:ARG5}, modifiers, named entities, etc.).
    \item The Konkani source sentence.
    \item A request for the Penman-notation AMR graph only, without explanation.
\end{enumerate}

The AMR graphs produced by Gemini use standard English concept nodes (e.g., \texttt{approve-01}, \texttt{increase-01}), consistent with the PropBank sense inventory, while the source sentence remains in Konkani. This is coherent with the AMR philosophy of language-neutral meaning representation.

\subsection{Human Verification}

All 1,100 automatically generated annotations were reviewed by a human expert with knowledge of both Konkani and AMR. The reviewer assessed each annotation for:
\begin{itemize}
    \item \textbf{Semantic faithfulness}: does the graph correctly represent the meaning of the Konkani sentence?
    \item \textbf{Structural correctness}: is the Penman notation syntactically valid (balanced parentheses, proper variable naming, valid relations)?
    \item \textbf{PropBank alignment}: are verb senses (e.g., \texttt{approve-01}) correctly selected from the PropBank frame inventory?
\end{itemize}

Annotations were accepted as-is, corrected, or rejected. The human reviewer accepted \textbf{97\%} of the automatically generated annotations without modification, with only 3\% requiring correction. No annotations were entirely rejected. This high acceptance rate provides strong evidence that Gemini 2.5 Pro produces reliable AMR graphs for Konkani sentences when provided with appropriate prompting.

\subsection{Dataset Statistics and Quality}

% tables/dataset_stats.tex
\begin{table}[t]
\centering
\small
\begin{tabular}{lr}
\toprule
\textbf{Statistic} & \textbf{Value} \\
\midrule
\multicolumn{2}{l}{\textit{Corpus size}} \\
\quad Total sentence--graph pairs & 1,100 \\
\quad Training set & 880 \\
\quad Validation set & 55 \\
\quad Test set & 165 \\
\midrule
\multicolumn{2}{l}{\textit{Sources}} \\
\quad AI4Bharat BPCC (random sample) & 1,000 \\
\quad Konkani newspaper & 100 \\
\midrule
\multicolumn{2}{l}{\textit{Sentence statistics}} \\
\quad Avg.\ sentence length (tokens) & 13.4 \\
\quad Min / Max sentence length & 4 / 36 \\
\midrule
\multicolumn{2}{l}{\textit{AMR graph statistics}} \\
\quad Avg.\ nodes per graph & 10.8 \\
\quad Avg.\ edges per graph & 13.8 \\
\midrule
\multicolumn{2}{l}{\textit{Annotation}} \\
\quad Annotator & Gemini 2.5 Pro \\
\quad Human verification acceptance rate & 97\% \\
\bottomrule
\end{tabular}
\caption{KonkaniAMR dataset statistics.}
\label{tab:dataset_stats}
\end{table}


Table~\ref{tab:dataset_stats} summarises the key statistics of KonkaniAMR. The dataset spans 1,100 sentences with an average length of 11.3 tokens. AMR graphs have an average of 8.4 nodes and 7.6 edges per graph.

\textbf{Qualitative examples.} We present three representative examples to illustrate the quality and diversity of the annotations.

\textbf{Example 1:} A straightforward administrative sentence with a named entity and a monetary quantifier:

\begin{quote}
\textit{Konkani:} मुरगाव पालिकेच्या 239 लाखांच्या शिलकी अर्थसंकल्पाक मान्यताय \\
\textit{(Translation:} The surplus budget of 239 lakhs of the Murgaon municipality was approved.\textit{)}
\end{quote}

\begin{verbatim}
(a / approve-01
   :ARG1 (b / budget
      :poss (g / government-organization
         :name (n / name :op1 "Murgaon"))
      :topic (s / surplus
         :quant (m / monetary-quantity
            :quant 23900000))))
\end{verbatim}

The graph correctly identifies \texttt{approve-01} as the root predicate, slots the budget as the entity being approved (\texttt{:ARG1}), links it to the Murgaon municipality via \texttt{:poss}, and encodes the monetary quantity in base units.

\textbf{Example 2:} A scalar change event with a precise quantity and measurement unit:

\begin{quote}
\textit{Konkani:} डिझेलाच्या दरांत प्रति लिटर 80 पयशांनी वाड जाल्या \\
\textit{(Translation:} Diesel prices have increased by 80 paise per litre.\textit{)}
\end{quote}

\begin{verbatim}
(i / increase-01
   :ARG1 (p / price
      :poss (d / diesel))
   :ARG2 (m / monetary-quantity
      :quant 80
      :unit (p2 / paise)
      :per (l / liter)))
\end{verbatim}

Here \texttt{increase-01} takes the diesel price as its theme (\texttt{:ARG1}) and the magnitude of change as \texttt{:ARG2}, with proper unit decomposition.

\textbf{Example 3:} An abstract political event capturing aspect and modality:

\begin{quote}
\textit{Konkani:} राजकी परिस्थिती खिणाखिणाक बदलत आसा \\
\textit{(Translation:} The political situation is changing moment by moment.\textit{)}
\end{quote}

\begin{verbatim}
(c / change-01
   :ARG1 (s / situation
      :mod (p / political))
   :manner (c2 / constant))
\end{verbatim}

The graph captures the continuously changing nature through the \texttt{:manner} relation with \texttt{constant}, correctly interpreting the Konkani aspectual construction ``खिणाखिणाक'' (moment-by-moment) as a manner modifier.

\subsection{Pretraining Corpus}

In addition to the annotated AMR dataset, we compiled a larger Konkani-language corpus for domain-adaptive pre-training. Using the same AI4Bharat BPCC dataset, we extracted all available Konkani target-side sentences from the corpus, processed with the \texttt{create\_pretraining.py} script. After deduplication, the corpus contains \textbf{189,332 unique Konkani sentences}. This corpus is used exclusively as unlabelled text for continued language pre-training (Section~\ref{sec:methodology}) and is not part of the AMR annotation pipeline.
